\documentclass[11pt,a4paper]{article}

% Compatible con pdfLaTeX, motor predeterminado de Overleaf.
\usepackage[utf8]{inputenc}
\usepackage[T1]{fontenc}
\usepackage[spanish,es-nodecimaldot]{babel}
\usepackage{geometry}
\usepackage{graphicx}
\usepackage{xcolor}
\usepackage{hyperref}
\usepackage{url}
\usepackage{booktabs}
\usepackage{tabularx}
\usepackage{tocloft}
\usepackage{array}
\usepackage{enumitem}
\usepackage{float}
\usepackage{listings}
\usepackage{fancyhdr}

\geometry{margin=2.4cm,headheight=14pt}
\setlength{\parskip}{0.45em}
\setlength{\parindent}{0pt}
\setlength{\emergencystretch}{3em}
\setlist[itemize]{leftmargin=1.5em,itemsep=0.2em,topsep=0.3em}
\setlist[enumerate]{leftmargin=1.7em,itemsep=0.25em,topsep=0.3em}
\renewcommand{\arraystretch}{1.2}
\setlength{\cftsecnumwidth}{3.0em}
\setlength{\cftsubsecnumwidth}{3.8em}
\urlstyle{same}
\Urlmuskip=0mu plus 2mu\relax

\definecolor{azul}{HTML}{244A63}
\definecolor{gris}{HTML}{F1F3F4}
\definecolor{gristexto}{HTML}{4C5661}
\definecolor{verde}{HTML}{2D6657}

\hypersetup{
  colorlinks=true,
  linkcolor=azul,
  urlcolor=verde,
  citecolor=azul,
  pdftitle={Fruta a la sartén: Memoria del proyecto final de SREX},
  pdfauthor={Irene Medina García y Vicente Rivas Monferrer}
}

\pagestyle{fancy}
\fancyhf{}
\lhead{\small Fruta a la sartén}
\rhead{\small Memoria del proyecto}
\cfoot{\thepage}
\renewcommand{\headrulewidth}{0.4pt}

\lstdefinestyle{csharp}{
  language={[Sharp]C},
  basicstyle=\ttfamily\footnotesize,
  keywordstyle=\color{azul}\bfseries,
  commentstyle=\color{gristexto}\itshape,
  stringstyle=\color{verde},
  numbers=left,
  numberstyle=\tiny\color{gristexto},
  numbersep=8pt,
  stepnumber=1,
  showstringspaces=false,
  breaklines=true,
  breakatwhitespace=false,
  tabsize=4,
  keepspaces=true,
  columns=fullflexible,
  frame=single,
  rulecolor=\color{gristexto},
  backgroundcolor=\color{gris},
  literate=
    {á}{{\'a}}1 {é}{{\'e}}1 {í}{{\'i}}1 {ó}{{\'o}}1 {ú}{{\'u}}1
    {Á}{{\'A}}1 {É}{{\'E}}1 {Í}{{\'I}}1 {Ó}{{\'O}}1 {Ú}{{\'U}}1
    {ñ}{{\~n}}1 {Ñ}{{\~N}}1 {ü}{{\"u}}1 {Ü}{{\"U}}1
    {¿}{{\textquestiondown}}1 {¡}{{\textexclamdown}}1
}

% Si existe la captura en la ruta indicada, se incluye automáticamente.
% Si no existe, el recuadro pendiente mantiene compilable el documento.
\newcommand{\captura}[3]{%
  \begin{figure}[H]
    \centering
    \IfFileExists{figuras/#1}{%
      \includegraphics[width=0.96\linewidth,height=0.48\textheight,keepaspectratio]
        {figuras/#1}%
    }{%
      \fcolorbox{azul}{gris}{%
        \parbox[c][4.1cm][c]{0.88\linewidth}{%
          \centering
          {\Large\textbf{CAPTURA PENDIENTE}}\\[0.6em]
          #3\\[0.8em]
          {\small Guardar como \path{figuras/#1}}%
        }%
      }%
    }%
    \caption{#2}
    \label{fig:#1}
  \end{figure}
}

\newcommand{\incluirscript}[3]{%
  \IfFileExists{Proyecto/Proyecto/Assets/_Proyecto/Scripts/#1}{%
    \lstinputlisting[style=csharp,caption={#2},label={#3}]
      {Proyecto/Proyecto/Assets/_Proyecto/Scripts/#1}%
  }{%
    \IfFileExists{Assets/_Proyecto/Scripts/#1}{%
      \lstinputlisting[style=csharp,caption={#2},label={#3}]
        {Assets/_Proyecto/Scripts/#1}%
    }{%
      \IfFileExists{scripts/#1}{%
        \lstinputlisting[style=csharp,caption={#2},label={#3}]
          {scripts/#1}%
      }{%
        \fcolorbox{azul}{gris}{\parbox{0.9\linewidth}{%
          \textbf{Código pendiente de adjuntar:} subir el archivo
          \texttt{#1} para que aparezca en este apéndice.}}%
      }%
    }%
  }%
}

\begin{document}

\hypersetup{pageanchor=false}
\begin{titlepage}
  \centering
  \vspace*{2cm}
  {\Large Sistemas de Realidad Extendida\par}
  \vspace{2cm}
  {\Huge\bfseries FRUTA A LA SARTÉN\par}
  \vspace{0.7cm}
  {\LARGE Memoria del proyecto final\par}
  \vfill
  {\Large Irene Medina García\\[0.3em]Vicente Rivas Monferrer\par}
  \vspace{1cm}
  {\large \today\par}
  \vfill
  \begin{minipage}{0.9\textwidth}
    \small
    Los datos en \textbf{negrita} señalan aspectos por completar, comprobar o
    documentar. La memoria distingue la configuración guardada en los archivos
    del proyecto de las pruebas realizadas en Meta XR Simulator y de las
    limitaciones que siguen pendientes.
  \end{minipage}
\end{titlepage}
\hypersetup{pageanchor=true}
\pagenumbering{arabic}

\tableofcontents
\newpage

\section{Descripción del proyecto}

\subsection{Temática y objetivo}

El proyecto propone una experiencia breve de realidad virtual ambientada en una
vivienda. El usuario recorre la casa, recoge piezas de fruta y las lleva a una
sartén para completar una receta. La cocina es el espacio principal de
interacción, mientras que las demás habitaciones forman parte del entorno
explorable.

El objetivo es integrar en una misma experiencia la configuración XR, la
interacción, la locomoción, el sonido espacial y una interfaz espacial. La
aplicación se plantea como realidad virtual para Meta Quest; no se ha planteado
como experiencia de realidad mixta.

\subsection{Experiencia prevista}

El usuario puede desplazarse por la casa, elegir un modo de locomoción, agarrar
fruta y depositarla en la zona receptora de la sartén. El marcador de la
interfaz muestra el progreso de la receta: dos plátanos, dos naranjas y una
piña. Se incluyen sonidos ambientales y efectos asociados a elementos y
acciones de la escena.

\textbf{Alcance de lo verificado.} La escena activa es
\path{Assets/_Proyecto/Scenes/Casa.unity}. Unity 6 importó la escena y el
\texttt{AudioMixer.mixer} en una ejecución batch que terminó correctamente. Esto
confirma la lectura de los assets, pero no el funcionamiento en tiempo de
ejecución. El intento actual de iniciar Meta XR Simulator se bloquea al crear el
swapchain (\texttt{XR\_ERROR\_RUNTIME\_FAILURE}, sin tipos de memoria
compatibles), por lo que no se afirma aquí una validación de la versión actual
en el simulador. Los scripts implementan el recuento de fruta y la detección de
la entrega en la sartén. Siguen pendientes las capturas y la prueba de ejecución
que pueda realizarse cuando el simulador arranque.

\section{Configuración del entorno}

\subsection{Proyecto y escena}

El proyecto utiliza Unity 6.0.83f1, Meta XR All-in-One SDK 209.0.0, Meta
OpenXR 2.7.0, URP 17.0.4 y AI Navigation 2.0.14. La escena principal de trabajo
es \path{Casa.unity}. El entorno se basa en el recurso de interiores low-poly
indicado en la sección de recursos externos.

La escena contiene elementos del ecosistema Meta XR para la cámara y la
interacción. El Camera Rig serializa el origen de seguimiento \texttt{Floor
Level}. La posición inicial y los ajustes XR se pueden inspeccionar en los
assets. La importación batch no valida su comportamiento durante la ejecución;
la comprobación actual en Meta XR Simulator sigue bloqueada por el error de
swapchain descrito anteriormente. No se presenta como una prueba en un visor
físico.

\captura{configuracion_xr.png}
  {Escena Casa y configuración del proyecto XR}
  {Capturar la escena y el Inspector del Camera Rig con el origen Floor Level.
  La imagen documenta la configuración serializada; no sustituye la prueba de
  ejecución pendiente.}

\subsection{Colliders de las paredes: resolución del vano}

Durante la preparación del escenario apareció un problema de colisión: aunque
una pared tenía un vano visible, el collider impedía que el usuario pasara por
él. Un collider simplificado que cubría la pared completa trataba como sólido
también el espacio vacío de la abertura.

La solución adoptada fue ajustar la colisión por elementos de la instancia de
la casa, añadiendo colliders a las partes sólidas de la pared y dejando libre el
volumen del vano. Así, la geometría de colisión representa los tramos que
realmente bloquean el paso, en vez de cerrar la abertura con una única caja. La
escena serializada contiene numerosos \texttt{BoxCollider} añadidos a elementos
del FBX; la importación del modelo no genera colliders automáticamente.

La escena serializa colliders en los tramos sólidos de la pared y deja sin
collider el volumen del vano, según la configuración de los objetos. La prueba
de atravesar la abertura y chocar con los tramos sólidos queda pendiente de
poder iniciar el simulador con la configuración actual.

\captura{colliders_paredes.png}
  {Colliders de pared y vano}
  {Mostrar el vano con los colliders visibles en el Editor: el hueco debe quedar
  libre y las partes sólidas deben conservar su colisión.}

\subsection{Superficies de suelo y teletransporte}

El suelo original estaba compuesto por muchos planos de \(1\times1\). Para
simplificar la definición de las superficies de destino del teletransporte, se
sustituyeron por cinco planos propios distribuidos por estancia:

\begin{itemize}
  \item \texttt{Suelo\_Cocina}
  \item \texttt{Suelo\_Dormitorio}
  \item \texttt{Suelo\_Salon}
  \item \texttt{Suelo\_Pasillo}
  \item \texttt{Suelo\_Aseo}
\end{itemize}

En la escena, estos objetos cuentan con \texttt{MeshFilter},
\texttt{NavMeshSurface} y \texttt{BoxCollider}. La división por habitaciones
simplifica la definición de las áreas de suelo y reduce la fragmentación
respecto al mosaico inicial. No se han registrado mediciones de rendimiento,
por lo que no se atribuye una mejora cuantitativa. La distribución está
serializada en la escena; su uso y accesibilidad deben confirmarse en ejecución.

\captura{suelos_teletransporte.png}
  {Los cinco planos de suelo}
  {Mostrar sus nombres en la jerarquía y los componentes relevantes en el
  Inspector, junto a las áreas de teletransporte.}

\subsection{Organización}

Los recursos propios se agrupan bajo \path{Assets/_Proyecto}, separados por
tipo de contenido. Esta organización permite distinguir los assets de la
experiencia de los paquetes y componentes instalados por Unity y Meta XR.

\begin{verbatim}
Assets/
|-- _Proyecto/
|   |-- Audio/
|   |-- Materials/
|   |-- Models/
|   |   |-- DearMousey/Stylized Fruits Props/
|   |   |-- LowPolyInteriorsFree/
|   |   `-- sarten.fbx
|   |-- Prefabs/Frutas/
|   |-- Scenes/Casa.unity
|   `-- Scripts/
|-- CompositionLayers/
|-- Oculus/
|-- Plugins/
|-- Resources/
|-- Screenshots/
|-- Settings/
|-- StreamingAssets/
|-- TextMesh Pro/
|-- TutorialInfo/
|-- UnityMCP/
`-- XR/
\end{verbatim}

\begin{itemize}
  \item \path{Audio/} contiene los clips y \path{AudioMixer.mixer}.
  \item \path{Materials/} contiene materiales y texturas del proyecto.
  \item \path{Models/} conserva los modelos y carpetas de procedencia.
  \item \path{Prefabs/Frutas/} agrupa las frutas interactuables reutilizables.
  \item \path{Scenes/} contiene la escena activa \path{Casa.unity}.
  \item \path{Scripts/} separa la lógica C\# propia de los componentes del SDK.
\end{itemize}

En el Hierarchy, \texttt{MUNDO}, \texttt{MUSICA}, \texttt{CONTROLLERS} e
\texttt{INTERFACES} son contenedores funcionales. Se escriben en mayúsculas
como convención visual para que los roots se distingan de los objetos de
contenido; las mayúsculas no tienen un efecto funcional en Unity. \texttt{MUNDO}
agrupa el escenario; \texttt{MUSICA}, el ambiente; \texttt{CONTROLLERS}, los
controladores de la aplicación; e \texttt{INTERFACES}, los paneles y elementos
de UI.

El Camera Rig y los Building Blocks de teletransporte y audio aprovechan
configuraciones reutilizables del SDK en vez de reconstruir manualmente cada
subsistema. \texttt{SartenController} permanece en \texttt{SartenSensor}, junto
al collider trigger, porque \texttt{OnTriggerEnter} debe recibir la entrada en
el objeto asociado al collider. No se agrupa bajo \texttt{CONTROLLERS} por esa
dependencia física.

\textbf{La escena conserva también \texttt{SUELO\_NAV}, que en los archivos
revisados solo presenta Transform y no una referencia funcional identificada;
no se le atribuye una función operativa.}

\captura{jerarquia_proyecto.png}
  {Jerarquía de Casa y carpetas principales}
  {Capturar el Hierarchy mostrando los roots funcionales, Camera Rig, sensores,
  Building Blocks y los suelos; acompañar con el árbol de carpetas del proyecto.}

\section{Interacción}

\subsection{Recogida y entrega de fruta}

La interacción se apoya en elementos del Meta Interaction SDK y en prefabs de
fruta. La lógica de entrega está implementada en \texttt{SartenController.cs},
que se encuentra en el objeto \texttt{SartenSensor}, asociado al collider
trigger de la sartén.

Cuando una fruta entra en el trigger, el script comprueba su etiqueta
(\texttt{Platano}, \texttt{Naranja} o \texttt{Pina}) y solicita registrar el
ingrediente en \texttt{FruitProgressController.cs}. Si todavía falta esa fruta
para completar la receta, se actualiza el marcador, se reproduce el efecto de
fritura y se destruye el objeto entregado. Un conjunto de identificadores evita
procesar dos veces la misma fruta.

Las cantidades objetivo configuradas en el script son dos plátanos, dos
naranjas y una piña. El progreso se presenta mediante contadores e iconos de
estado en la interfaz.

\captura{interaccion_fruta.png}
  {Recogida y entrega de fruta}
  {Capturar al usuario recogiendo una fruta y depositándola en la sartén durante
  la ejecución, con el marcador visible.}

\subsection{Interacción cercana y a distancia}

La fruta constituye la interacción cercana de la experiencia. Los prefabs
incluyen componentes del Interaction SDK para agarre de mano y una rutina de
instalación de agarre a distancia. La interfaz ofrece interacción a distancia
mediante los componentes de Canvas y el rayo del usuario; el poke proporciona
una modalidad cercana para los controles.

Los prefabs y componentes de interacción están serializados en la escena y en
los prefabs de fruta. La recogida, la entrega y la interacción por rayo o manos
deben volver a confirmarse en ejecución; el intento actual de iniciar el
simulador no llega a esa fase. La captura pendiente servirá como evidencia cuando
se pueda realizar la prueba.

\captura{rayo_poke.png}
  {Interacción del usuario con la fruta y la interfaz}
  {Aportar evidencia del agarre y de la selección de un control con el rayo; si
  se usa poke, mostrar también esa interacción cercana.}

\subsection{Estado de la función temática}

El control de la interfaz asociado al marcador no reinicia actualmente la
receta. En \texttt{ProgressRefreshButton.cs}, la acción llama a
\texttt{RefreshProgressUI()}, que vuelve a dibujar el estado existente, pero no
pone a cero los contadores ni repone las frutas. Por tanto, no debe describirse
como un botón de reinicio funcional hasta que se implemente y se pruebe esa
lógica.

\captura{trigger_sarten.png}
  {Trigger de la sartén y controlador de entrega}
  {Mostrar el collider trigger y el componente \texttt{SartenController} en el
  Inspector, junto al sensor de la sartén.}

\section{Locomoción}

\subsection{Modos de desplazamiento}

El script \texttt{LocomotionModeController.cs} conecta el desplegable de la
interfaz con \texttt{MovingSetting} del Meta Interaction SDK. Permite habilitar
desplazamiento continuo (\textit{Slide}), teletransporte, ambos modos o
ninguno, y asigna el estilo correspondiente cuando se selecciona uno de los dos
modos individuales.

Para activar o desactivar los objetos asociados a cada modo, el script consulta
mediante reflexión campos internos del SDK. Depende de nombres internos de sus
componentes, por lo que habría que revisarla si se actualiza el SDK. La
alternancia en ejecución queda pendiente de confirmación con la versión concreta
del paquete cuando el simulador pueda arrancar.

\subsection{Teletransporte y confort}

La escena contiene tres bloques \texttt{Teleport Interactable} y un objeto
llamado explícitamente \texttt{Teleport Hotspot}. Esto no basta para demostrar
por sí solo cómo se distribuyen las zonas de destino. La cantidad, ubicación y
accesibilidad deben comprobarse durante la ejecución.

En la instancia del Meta XR Rig se serializa el nivel de confort
\texttt{High} (\path{_comfortLevel = 3}). Su efecto debe confirmarse durante la
ejecución. No se afirma que el giro brusco (\textit{Snap Turn}) esté configurado,
al no aparecer en los archivos revisados.

\captura{locomocion.png}
  {Selector de locomoción}
  {Capturar el desplegable mostrando Slide, teletransporte, ambos y ninguno, y
  evidenciar los modos durante la ejecución.}

\captura{teletransporte_confort.png}
  {Áreas de teletransporte y ajuste de confort}
  {Mostrar la distribución real de los destinos en la casa y el ajuste High del
  rig; añadir una prueba de desplazamiento.}

\section{Sonido 3D}

\subsection{Fuentes presentes}

La escena incluye una fuente ambiental en bucle y emisores asociados a
habitaciones y objetos, entre ellos \texttt{SONIDO\_WC},
\texttt{SONIDO\_SARTEN} y \texttt{SONIDO\_FREIR}. El script
\texttt{SartenController.cs} reproduce el clip de fritura cuando se registra
correctamente una fruta. El volumen ambiental se controla mediante
\texttt{AmbientVolumeController.cs}; los efectos espaciales, mediante
\texttt{SpatialEffectsController.cs}.

La escena serializa tres Building Blocks de Spatial Audio y las fuentes
espaciales incorporan componentes Meta XR Audio con espacialización y acústica
habilitadas. Esto describe la configuración de los assets, no una prueba auditiva
actual; los valores acústicos de sala se presentan como aproximaciones, no como
mediciones.

La escena también contiene configuración acústica de sala mediante el modelo
\textit{Shoebox}, con dimensiones aproximadas de \(13\times4\times14\) m y
\texttt{clutterFactor = 0.5}. Son valores de aproximación al escenario; no se
dispone de una validación acústica medida. La opción serializada de difracción
no demuestra por sí sola difracción geométrica activa.

\captura{audio_fuentes.png}
  {Fuentes ambientales y espacializadas}
  {Mostrar la posición de los emisores y los ajustes de atenuación en el
  Inspector; incluir una evidencia auditiva o de prueba en ejecución.}

\subsection{Audio Mixer y enrutamiento}

\path{AudioMixer.mixer} contiene los grupos \texttt{Master}, \texttt{Ambiente} y
\texttt{Efectos}, con un snapshot inicial a \(0\) dB para los tres. En
\path{Casa.unity}, la fuente ambiental está enrutada a \texttt{Ambiente} y los
tres emisores de Spatial Audio a \texttt{Efectos}. La importación batch de Unity
terminó correctamente tras estos cambios. Falta comprobar visualmente el Mixer
en el Editor y probar durante la ejecución que los controles de la interfaz
afectan al audio esperado; la lectura correcta de los assets no demuestra por sí
sola ese comportamiento.

\captura{audio_mixer.png}
  {Audio Mixer}
  {Mostrar los grupos Master, Ambiente y Efectos, sus niveles iniciales y las
  fuentes enrutadas a cada grupo.}

\captura{acustica_sala.png}
  {Configuración acústica y dimensiones del espacio}
  {Mostrar el modelo Shoebox, las dimensiones 13 x 4 x 14 m y clutterFactor 0.5.}

\subsection{Procedencia de los audios}

En \path{usados.txt} se anotan referencias de Pixabay. Se comprobaron los títulos,
autores y licencias indicados en sus páginas:

\begin{itemize}
  \item \textit{Wet Squelch}, autor Wolff-SFX:
        \href{https://pixabay.com/sound-effects/film-special-effects-wet-squelch-425547/}{ficha del sonido}.
        La página indica la Pixabay Content License.
  \item \textit{Room Tone 1}, autor stevebfx (Steven Blasini):
        \href{https://pixabay.com/sound-effects/film-special-effects-room-tone-1-579801/}{ficha del sonido}.
        La página indica la Pixabay Content License.
  \item \textit{The WC sound}, autor Romano.Tercero (Freesound), publicado en
        Pixabay:
        \href{https://pixabay.com/es/sound-effects/familiar-the-wc-sound-66433/}{ficha del sonido}.
        La página indica la Pixabay Content License.
\end{itemize}

\textbf{Queda pendiente confirmar qué archivo local (\texttt{ambient.mp3},
\texttt{freir.mp3}, \texttt{sarten.mp3} o \texttt{wc.mp3}) corresponde a cada
referencia y qué fuente de la escena lo reproduce.} No se atribuye por nombre de
archivo una correspondencia que no se haya verificado escuchando o revisando el
clip.

\section{Interfaces}

La escena contiene una interfaz espacial con controles para el volumen ambiental,
el modo de locomoción y los efectos espaciales. El volumen se conecta al
slider mediante \texttt{AmbientVolumeController.cs}; la locomoción, al
desplegable mediante \texttt{LocomotionModeController.cs}; y los efectos
espaciales, al toggle asociado mediante \texttt{SpatialEffectsController.cs}.

La interfaz incluye además un control de actualización del marcador. Tal como
está implementado, este solo actualiza la representación del progreso: no
reinicia la receta. Por tanto, no debe contarse como la función temática de
reinicio solicitada hasta que se corrija.

El Canvas está en modo World Space. La escena contiene componentes de
interacción Ray y Poke para los paneles. El uso de la interfaz debe confirmarse
en ejecución cuando el simulador pueda arrancar. La captura pendiente documenta
el Canvas y sus componentes; la presencia de Ray/Poke no se presenta por sí sola
como prueba del uso de prefabs concretos del Meta UI Set.

\captura{interfaz_espacial.png}
  {Canvas espacial y controles de la aplicación}
  {Mostrar el modo World Space, la ubicación del panel, sus componentes o
  prefabs Meta UI Set y la interacción por rayo.}

\section{Recursos externos}

\begin{table}[H]
\centering
\small
\begin{tabularx}{\textwidth}{>{\raggedright\arraybackslash}p{0.20\textwidth}
                            X}
\toprule
\textbf{Recurso} & \textbf{Procedencia y comprobaciones}\\
\midrule
Escenario e interiores & \textit{DOLLHOUSE -- FREE -- Modular Low Poly
Interiors: Furniture \& Room Kit}, Unity Asset Store, producto 411314.
\href{https://assetstore.unity.com/packages/3d/props/interior/dollhouse-free-modular-low-poly-interiors-furniture-room-kit-411314}{Ficha del paquete}.\\
Frutas & \textit{Stylized Fruit Props [FREE]}, Unity Asset Store, producto
391728. \href{https://assetstore.unity.com/packages/3d/props/stylized-fruit-props-free-391728}{Ficha del paquete}.\\
Sartén & \textit{Frying pan lowpoly}, modelo de Sketchfab de rombikgen.
\href{https://sketchfab.com/3d-models/frying-pan-lowpoly-89432a8085ff409bb6b00b6db4e20a29}{Ficha del modelo}.
Licencia Creative Commons Atribución 4.0 Internacional (CC BY 4.0). La
atribución debe conservar el título, el autor y el enlace al modelo.\\
Efectos de sonido & Las páginas de Pixabay y sus autores/licencias se identifican
en el apartado de sonido. Falta asociar de forma inequívoca cada enlace con cada
archivo de audio local.\\
\bottomrule
\end{tabularx}
\caption{Recursos externos utilizados y comprobaciones pendientes}
\end{table}

\section{Organización del trabajo}

El reparto debe reflejar el trabajo realmente realizado por cada integrante.
Los nombres del grupo son Irene Medina García y Vicente Rivas Monferrer; las
tareas concretas se dejan pendientes para no atribuir responsabilidades sin
confirmación.

\begin{table}[H]
\centering
\begin{tabularx}{\textwidth}{>{\raggedright\arraybackslash}p{0.32\textwidth}X}
\toprule
\textbf{Integrante} & \textbf{Tareas realizadas y responsabilidad}\\
\midrule
Irene Medina García & \textbf{[Completar con las tareas efectivamente
realizadas y sus evidencias.]}\\
Vicente Rivas Monferrer & \textbf{[Completar con las tareas efectivamente
realizadas y sus evidencias.]}\\
\bottomrule
\end{tabularx}
\caption{Reparto de responsabilidades del grupo}
\end{table}

\section{Comprobaciones pendientes antes de la entrega}

\begin{enumerate}
  \item Implementar el reinicio real de la receta y, si corresponde, reponer las
        frutas. Esto requiere modificar C\#, que no se ha cambiado en esta
        revisión.
  \item Confirmar en el Editor los grupos y enrutamientos del Audio Mixer, y
        probar los controles y la atenuación en ejecución.
  \item Resolver el bloqueo de Meta XR Simulator (\texttt{xrCreateSwapchain},
        sin tipos de memoria compatibles) y repetir las pruebas de interacción,
        locomoción, confort, audio e interfaz. No se han cambiado preferencias
        de GPU ni de Windows.
  \item Asociar cada archivo de audio con su página de procedencia y añadir la
        atribución correspondiente en la memoria.
  \item Completar el reparto real del trabajo y sustituir los marcadores de
        captura por imágenes de la escena y de las pruebas.
\end{enumerate}

\section{Conclusiones}

El proyecto cuenta con una escena de vivienda adaptada para la experiencia, una
configuración de colliders, cinco planos de suelo propios, scripts para registrar
la entrega de fruta y mostrar el progreso, y componentes de locomoción, sonido e
interfaz. El Audio Mixer ya contiene grupos diferenciados y las fuentes de la
escena están enrutadas a ellos. La importación batch de Unity terminó
correctamente, pero no equivale a una prueba de ejecución.

En la revisión actual, el inicio de Meta XR Simulator queda bloqueado antes de
probar la experiencia por un error de asignación de memoria al crear el
swapchain. Siguen pendientes el reinicio temático (requiere C\#), la validación
en ejecución, la correspondencia de los audios con sus fuentes, el reparto real
del trabajo y las capturas. No se afirma una prueba en un visor físico.
\appendix
\section{Código fuente de los scripts}

Los listados siguientes se leen directamente desde
\path{Assets/_Proyecto/Scripts}. De este modo, al compilar la memoria se incluye
el código presente en el proyecto y se evita mantener una copia manual que
pueda quedar desactualizada. El comando busca los archivos tanto en la carpeta
del proyecto Unity como en \path{scripts/}; si se compila en Overleaf sin esos
archivos, aparecerá un aviso en su lugar.

\subsection{\texttt{AmbientVolumeController.cs}}
\incluirscript{AmbientVolumeController.cs}{Control del volumen ambiental}{lst:ambient}

\subsection{\texttt{ExitDialogController.cs}}
\incluirscript{ExitDialogController.cs}{Diálogo de salida y estado de la tarea}{lst:exit}

\subsection{\texttt{FruitProgressController.cs}}
\incluirscript{FruitProgressController.cs}{Registro y presentación del progreso de fruta}{lst:progress}

\subsection{\texttt{HelpDialogController.cs}}
\incluirscript{HelpDialogController.cs}{Visibilidad del panel de ayuda}{lst:help}

\subsection{\texttt{LocomotionModeController.cs}}
\incluirscript{LocomotionModeController.cs}{Selección del modo de locomoción}{lst:locomotion}

\subsection{\texttt{ProgressRefreshButton.cs}}
\incluirscript{ProgressRefreshButton.cs}{Actualización visual del progreso}{lst:refresh}

\subsection{\texttt{SartenController.cs}}
\incluirscript{SartenController.cs}{Detección de fruta en el sensor de la sartén}{lst:pan}

\subsection{\texttt{SpatialEffectsController.cs}}
\incluirscript{SpatialEffectsController.cs}{Silenciado y activación de efectos espaciales}{lst:spatial}

\end{document}
